\documentclass[11pt,reqno]{amsart}

\usepackage[utf8]{inputenc}
\usepackage[T1]{fontenc}
\usepackage{amsmath,amssymb,amsthm}
\usepackage{lmodern}
\usepackage{mathtools}
\usepackage{mathrsfs}
\usepackage[shortlabels]{enumitem}
\usepackage[margin=1.15in]{geometry}
\usepackage[protrusion=true,expansion=false]{microtype}
\usepackage{booktabs}
\usepackage{tikz-cd}
\usepackage[colorlinks=true,linkcolor=blue!70!black,citecolor=green!50!black,urlcolor=blue!80!black]{hyperref}

% Theorem Environments
\theoremstyle{plain}
\newtheorem{theorem}{Theorem}[section]
\newtheorem{lemma}[theorem]{Lemma}
\newtheorem{proposition}[theorem]{Proposition}
\newtheorem{corollary}[theorem]{Corollary}

\theoremstyle{definition}
\newtheorem{definition}[theorem]{Definition}
\newtheorem{remark}[theorem]{Remark}

\numberwithin{equation}{section}

% Math Operators
\DeclareMathOperator{\Spec}{Spec}
\DeclareMathOperator{\Supp}{Supp}
\DeclareMathOperator{\Ann}{Ann}
\DeclareMathOperator{\Hom}{Hom}
\DeclareMathOperator{\codim}{codim}
\DeclareMathOperator{\Sing}{Sing}
\DeclareMathOperator{\Reg}{Reg}
\DeclareMathOperator{\NonNormal}{NonNormal}
\DeclareMathOperator{\Miss}{Miss}
\DeclareMathOperator{\Frac}{Frac}
\DeclareMathOperator{\Norm}{Norm}

\newcommand{\A}{\mathbb{A}}
\newcommand{\C}{\mathbb{C}}
\newcommand{\R}{\mathbb{R}}
\newcommand{\Z}{\mathbb{Z}}
\newcommand{\N}{\mathbb{N}}
\newcommand{\Q}{\mathbb{Q}}
\newcommand{\cO}{\mathcal{O}}
\newcommand{\cF}{\mathcal{F}}
\newcommand{\cQ}{\mathcal{Q}}
\newcommand{\fc}{\mathfrak{c}}
\newcommand{\fm}{\mathfrak{m}}
\newcommand{\fp}{\mathfrak{p}}

\title[Normal Loci of Non-Properness Sets of Keller Maps]{Normal Loci of Non-Properness Sets of Keller Maps\\and the Conductor Principle}

\author{}
\address{}
\date{August 2026}

\begin{document}

\begin{abstract}
Let $F: \A^n_\C \to \A^n_\C$ be a polynomial Keller map ($\det JF \in \C^\times$). Let $S_F \subset \A^n_\C$ denote its non-properness locus endowed with its reduced scheme structure, and let $S_F^{\mathrm{nor}} \subset S_F$ be its normal locus. We prove that the restricted morphism
\[
F^{-1}(S_F^{\mathrm{nor}}) \longrightarrow S_F^{\mathrm{nor}}
\]
is a finite, \'{e}tale, and surjective morphism. Consequently, the missing value locus $\Miss(F) = \A^n_\C \setminus F(\A^n_\C)$ is contained in the non-normal locus of $S_F$, which coincides with the vanishing locus of the conductor ideal of the normalization:
\[
\Miss(F) \subseteq \NonNormal(S_F) = V(\fc_{S_F^\nu / S_F}).
\]
In particular, if $S_F$ is normal (e.g.\ smooth), then $F$ is globally surjective, extending and refining Jelonek's 2022 smooth criterion.

The proof is self-contained modulo standard foundational results in algebraic geometry:
\begin{enumerate}[(i)]
\item Zariski's Main Theorem and the identification of relative normalization with function-field normalization for normal varieties (Stacks Tag 02LR, Tag 03GR);
\item the pure codimension-1 affine boundary theorem (Stacks Tag 0BCV);
\item Abhyankar's regular divisor theorem (Stacks Tag 0EYH / Lemma 58.31.6) to exclude interior-boundary collision over regular points;
\item Serre's $R_1+S_2$ criterion and height-one prime intersection algebraic Hartogs rigidity (Stacks Tag 0AVB) for affine schemes (Stacks Tag 01HX).
\end{enumerate}
\end{abstract}

\maketitle

\tableofcontents

\newpage

%=============================================================================
\section{Introduction and Main Results}
%=============================================================================

\subsection{Background}
Let $k = \C$ be the field of complex numbers. The classical Jacobian Conjecture of Keller \cite{Keller1939} asserts that any polynomial map $F: \A^n_\C \to \A^n_\C$ with non-vanishing constant Jacobian determinant $\det JF \in \C^\times$ (a \emph{Keller map}) is an invertible polynomial automorphism. By the Ax--Grothendieck theorem and a fundamental result of Bia{\l}ynicki-Birula and Rosenlicht \cite{BiaRosen1962}, every injective polynomial endomorphism of $\A^n_\C$ is an automorphism. In the context of non-proper polynomial mappings and asymptotic values, understanding the surjectivity and the locus of omitted values
\[
\Miss(F) := \A^n_\C \setminus F(\A^n_\C)
\]
plays a central role.

In 1993, Jelonek \cite{Jelonek1993} initiated the systematic geometric study of the non-properness locus $S_F$ of a dominant polynomial mapping $F: X \to Y$, defined as the set of points $y \in Y$ at which $F$ is not topologically or algebraically proper. Jelonek established that for any dominant polynomial map, $S_F$ is a pure codimension-1 algebraic subvariety (a hypersurface) of $Y$, or is empty. It is well known that any omitted value of $F$ must lie in $S_F$:
\[
\Miss(F) \subseteq S_F.
\]

In 2022, Jelonek \cite{Jelonek2022} proved that \emph{if the non-properness hypersurface $S_F$ is smooth, then $F$ is globally surjective} ($\Miss(F) = \emptyset$). 

The purpose of this paper is to uncover the precise geometric mechanism governing the non-properness locus of Keller maps and to upgrade Jelonek's smooth criterion to a pointwise structural theorem on the normal locus $S_F^{\mathrm{nor}}$.

\subsection{Main Theorems}
Throughout this paper, let $F: X = \A^n_\C \to Y = \A^n_\C$ be a Keller map. Let $S = S_F \subset Y$ denote its non-properness set endowed with its reduced subscheme structure. Let
\[
S_F^{\mathrm{nor}} := \left\{ y \in S_F : \cO_{S_F, y} \text{ is an integrally closed domain (normal local ring)} \right\}
\]
denote the normal locus of $S_F$.

Our first main result is the following structural rigidity theorem:

\begin{theorem}[Normal-Locus Finite \'{E}tale Theorem]\label{thm:normal-locus}
Let $F: \A^n_\C \to \A^n_\C$ be a polynomial Keller map ($\det JF \in \C^\times$). Then the restricted morphism
\[
F^{-1}(S_F^{\mathrm{nor}}) \longrightarrow S_F^{\mathrm{nor}}
\]
is a finite, \'{e}tale, and surjective morphism.
\end{theorem}

Theorem~\ref{thm:normal-locus} immediately establishes that missing values are mathematically forbidden from occurring at any normal point of $S_F$:

\begin{corollary}[The Conductor Principle]\label{cor:conductor}
Let $F: \A^n_\C \to \A^n_\C$ be a polynomial Keller map. Then:
\[
\Miss(F) \subseteq \NonNormal(S_F).
\]
Equivalently, if $\nu: S_F^\nu \to S_F$ is the normalization of $S_F$ and $\fc_{S_F^\nu / S_F} = \Ann_{\cO_{S_F}}(\nu_* \cO_{S_F^\nu} / \cO_{S_F})$ is the conductor ideal, then:
\[
\Miss(F) \subseteq V(\fc_{S_F^\nu / S_F}).
\]
\end{corollary}

\begin{corollary}[Global Normal Surjectivity]\label{cor:normal-surjectivity}
If $S_F$ is a globally normal variety (in particular, if $S_F$ is smooth), then $F$ is surjective:
\[
F(\A^n_\C) = \A^n_\C.
\]
\end{corollary}

\begin{corollary}[Defect Rigidity]\label{cor:defect-rigidity}
Let $W \subseteq S_F^{\mathrm{nor}}$ be any connected component. Then there exists an integer $r_W \ge 1$ such that
\[
\#F^{-1}(y) = r_W \quad \text{for all } y \in W.
\]
In particular, the fiber cardinality cannot drop or jump on the normal locus.
\end{corollary}

\subsection{The Role of Normality in Boundary Rigidity}

Normality enters the proof at two distinct stages.

First, let $Z \subset S_F^{\mathrm{nor}}$ be the normal irreducible affine neighborhood introduced in Section~3.  The Boundary Collision Lemma shows that every boundary point of the closure of $F^{-1}(Z)$ must map into $\Sing(Z)$.  Since $Z$ is normal, Serre's condition $R_1$ implies that $\Sing(Z)$ has codimension at least $2$.  Thus Abhyankar's regular-divisor theorem controls the codimension-$1$ geometry, while normality forces every remaining boundary phenomenon into higher codimension.

Second, since $F^{-1}(Z) \to Z$ is \'{e}tale and $Z$ is normal, $T := F^{-1}(Z)$ is itself normal.  Hence the normalization $\widetilde{C} \to C$ of its closure is an isomorphism over $T$. The normal scheme $\widetilde{C}$ satisfies the height-one intersection property
\[
\Gamma(\widetilde{C}_i, \cO_{\widetilde{C}_i}) = \bigcap_{\operatorname{ht}\fp=1} (\cO_{\widetilde{C}_i})_\fp
\]
on each irreducible component.  Since the complement $\widetilde{C}_i \setminus T_i$ has codimension at least $2$, restriction induces
\[
\Gamma(\widetilde{C}_i, \cO_{\widetilde{C}_i}) \cong \Gamma(T_i, \cO_{T_i}).
\]
Finally, $T_i$ is affine.  Therefore the equality of coordinate rings upgrades to an isomorphism of affine schemes $T_i \cong \widetilde{C}_i$.

Conceptually, the proof separates the boundary problem into two parts:
\[
\text{regular codimension-$1$ geometry} \quad\longrightarrow\quad \text{Abhyankar collision exclusion},
\]
and
\[
\text{residual codimension-$\ge 2$ geometry} \quad\longrightarrow\quad \text{Hartogs rigidity plus affineness}.
\]
This is the mechanism that allows the argument to work over the normal locus rather than requiring the entire non-properness set to be smooth.

\subsection{Relation to Prior Art}

The present argument is closely related to Jelonek's 2022 proof in the smooth case, but the mechanism used to control the boundary is different.

Both approaches pass to a finite completion supplied by Zariski's Main Theorem and reduce the problem to controlling the closure of the preimage of the non-properness set inside that finite completion.  In the smooth case, local structure over the smooth discriminant can be used to rule out boundary accumulation directly.

Here we separate the problem into codimension $1$ and higher codimension. Over the regular locus, Abhyankar's regular-divisor theorem rules out a collision between an interior divisor and a boundary divisor.  Consequently, all remaining boundary points lie over $\Sing(Z)$.  Since $Z$ is normal, Serre's condition $R_1$ places this singular locus in codimension at least $2$. The remaining boundary is then removed by the height-one Hartogs property of the normalization together with affineness.

Thus the argument does not require a global smoothness assumption on $S_F$. Instead, it proves the pointwise statement that
\[
F^{-1}(S_F^{\mathrm{nor}}) \longrightarrow S_F^{\mathrm{nor}}
\]
is finite, \'{e}tale, and surjective.

We also emphasize that this statement is different from the stronger normality direction considered in the withdrawn Jelonek--\.{Z}o{\l}\k{a}dek preprint \cite{JelonekZoladek2020}.  We do not claim that a normal irreducible non-properness set is impossible.  Our conclusion is instead that no normal point of $S_F$ can be an omitted value.

%=============================================================================
\section{The Canonical Finite Completion and \texorpdfstring{$S_F = \pi(L_\infty)$}{SF = pi(L)}}
%=============================================================================

\subsection{The Dominance of Keller Maps}
Let $F = (F_1, \dots, F_n): X = \A^n_\C \to Y = \A^n_\C$ be a polynomial map with $\det JF \in \C^\times$.

\begin{proposition}\label{prop:keller-etale}
$F$ is an \'{e}tale, quasi-finite, and dominant morphism.
\end{proposition}
\begin{proof}
Consider the coordinate ring presentation
\[
B = \C[x_1, \dots, x_n] \cong \C[y_1, \dots, y_n][x_1, \dots, x_n] / (F_1(x)-y_1, \dots, F_n(x)-y_n).
\]
The Jacobian matrix of the relations with respect to $x_j$ is precisely $JF = (\partial F_i / \partial x_j)$, whose determinant is an invertible constant in $\C^\times$. Thus, this is a standard smooth presentation of relative dimension 0. Since \'{e}tale is equivalent to smooth of relative dimension 0 (Stacks Tag \href{https://stacks.math.columbia.edu/tag/02GH}{02GH}), $F$ is \'{e}tale.

Every \'{e}tale morphism of schemes is an open map (Stacks Tag \href{https://stacks.math.columbia.edu/tag/02GH}{02GH}, Lemma 29.37.13). Since $X = \A^n_\C$ is non-empty, $F(X)$ is a non-empty Zariski open subset of the irreducible variety $Y = \A^n_\C$. In an irreducible variety, every non-empty open subset is dense. Hence $F$ is dominant, which implies that the pullback homomorphism
\[
F^*: \C[y_1, \dots, y_n] \hookrightarrow \C[x_1, \dots, x_n]
\]
is injective.
\end{proof}

\subsection{Identification of Relative Normalization and Field Normalization}
Let $K = K(X)$ be the function field of $X = \A^n_\C$.

\begin{lemma}[Canonical Finite Normalization Completion]\label{lem:finite-completion}
Let $F: X = \A^n_\C \to Y = \A^n_\C$ be a dominant quasi-finite morphism. Let $\pi: V \to Y$ be the normalization of $Y$ in the finite field extension $K(Y) \subset K(X)$.
Then there exists a canonical factorization
\begin{equation}\label{eq:finite-completion}
\begin{tikzcd}
X = \A^n_\C \arrow[r, hook, "j"] \arrow[rd, "F"'] & V \arrow[d, "\pi"] \\
& Y = \A^n_\C
\end{tikzcd}
\end{equation}
where $j: X \hookrightarrow V$ is a quasi-compact dense open immersion and $\pi: V \to Y$ is a finite morphism.
\end{lemma}
\begin{proof}
Let $Y^{\mathrm{rel}} \to Y$ be the relative normalization of $Y$ in the morphism $F: X \to Y$. Since $F$ is quasi-finite and separated, Zariski's Main Theorem (Stacks Tag \href{https://stacks.math.columbia.edu/tag/02LR}{02LR}, Lemma 37.43.2; EGA IV \cite{EGA4}, 8.12.6) yields a canonical factorization $X \xrightarrow{j} Y^{\mathrm{rel}} \to Y$ where $j$ is a quasi-compact open immersion. Since $Y = \A^n_\C$ is a Nagata scheme, $F$ is of finite type, and $X$ is reduced, the relative normalization $Y^{\mathrm{rel}} \to Y$ is finite (Stacks Tag \href{https://stacks.math.columbia.edu/tag/03GR}{03GR}, Lemma 29.54.15).

We now prove that $Y^{\mathrm{rel}} \cong V$. It suffices to check this over affine open subsets of $Y$.
Let $U = \Spec A \subset Y$ be any non-empty affine open subscheme, and let $W = F^{-1}(U)$. Since $F$ is dominant, $W \neq \emptyset$.
The relative normalization over $U$ is given by $(Y^{\mathrm{rel}})|_U = \Spec A_{\mathrm{rel}}$, where $A_{\mathrm{rel}}$ is the integral closure of $A$ in $\Gamma(W, \cO_X)$ (Stacks Tag \href{https://stacks.math.columbia.edu/tag/035F}{035F}, Lemma 29.54.1).

On the other hand, the field normalization over $U$ is $V|_U = \Spec \bar{A}$, where $\bar{A} = \{ c \in K : c \text{ is integral over } A\}$.
Since $W$ is a non-empty open subscheme of the integral scheme $X$, we have a canonical injection $\Gamma(W, \cO_X) \hookrightarrow K(X) = K$. Thus $A_{\mathrm{rel}} \subseteq \bar{A}$.

Conversely, let $c \in \bar{A}$. Then $c$ satisfies a monic equation $c^m + a_{m-1}c^{m-1} + \dots + a_0 = 0$ with $a_i \in A$.
For any point $x \in W$, the canonical map $A \to \cO_{X,x}$ shows that $c$ is integral over the local ring $\cO_{X,x}$.
Since $X = \A^n_\C$ is a normal variety, $\cO_{X,x}$ is an integrally closed domain with fraction field $\Frac \cO_{X,x} = K(X) = K$.
Therefore, $c \in \cO_{X,x}$ for every $x \in W$.
This implies $c$ is a regular function on $W$, so $c \in \Gamma(W, \cO_X)$. Since $c$ is integral over $A$, we obtain $c \in A_{\mathrm{rel}}$.

Thus $A_{\mathrm{rel}} = \bar{A}$, which proves $Y^{\mathrm{rel}} \cong V$. Therefore, $j: X \hookrightarrow V$ is a quasi-compact open immersion and $\pi: V \to Y$ is finite.
\end{proof}

We denote the \emph{boundary at infinity} by:
\[
L_\infty := V \setminus X.
\]

\subsection{Exact Identity \texorpdfstring{$S_F = \pi(L_\infty)$}{SF = pi(L)}}
We define the non-properness locus $S_F \subset Y$:
\begin{definition}\label{def:nonproperness}
The underlying topological set of the non-properness locus is defined as:
\[
|S_F| := \left\{ y \in Y : \nexists \text{ Zariski open } U \ni y \text{ such that } F^{-1}(U) \to U \text{ is proper} \right\}.
\]
\end{definition}

\begin{proposition}\label{prop:boundary-image}
Set-theoretically, $|S_F| = \pi(L_\infty)$. Consequently, $|S_F|$ is a closed algebraic subset of $Y$, which we endow with its reduced subscheme structure $S_F$. Furthermore, $\Miss(F) \subseteq S_F$.
\end{proposition}
\begin{proof}
($\subseteq$): Suppose $y \notin \pi(L_\infty)$. Since $\pi$ is finite, $\pi(L_\infty)$ is a closed subset of $Y$. Thus, there exists a Zariski open neighborhood $U \ni y$ such that $U \cap \pi(L_\infty) = \emptyset$. This implies $\pi^{-1}(U) \subset X$. Therefore, $F^{-1}(U) = \pi^{-1}(U) \to U$ is finite, hence proper. Thus $y \notin S_F$, which shows $S_F \subseteq \pi(L_\infty)$.

($\supseteq$): Suppose $y \in \pi(L_\infty)$, and assume for contradiction that $y \notin S_F$. Then there exists an affine open neighborhood $U = \Spec A \ni y$ such that $F^{-1}(U) \to U$ is proper. Since $F$ is quasi-finite, the restriction $F^{-1}(U) \to U$ is proper and quasi-finite, hence finite (Stacks Tag \href{https://stacks.math.columbia.edu/tag/02LS}{02LS}, Lemma 37.44.1). Thus $F^{-1}(U) = \Spec B$, where $B$ is a finite $A$-algebra.

Since $F$ is dominant, $F^{-1}(U)$ is a non-empty open subvariety of $X = \A^n_\C$. Hence $B$ is an integral domain, normal, and has quotient field $\Frac B = K(X) = K$. Let $\bar{A}$ denote the integral closure of $A$ in $K$ (so $\pi^{-1}(U) = \Spec \bar{A}$). Because $B/A$ is finite, $B \subseteq \bar{A}$. Conversely, let $c \in \bar{A}$. Then $c$ is integral over $A$, and hence integral over $B$. Since $B$ is normal (integrally closed in $K$), $c \in B$. Thus $B = \bar{A}$, which implies
\[
F^{-1}(U) = \pi^{-1}(U).
\]
However, $y \in \pi(L_\infty)$ implies that $\pi^{-1}(U)$ contains a point $a \in L_\infty = V \setminus X$, contradicting $F^{-1}(U) \subset X$. Thus $y \in S_F$, proving $\pi(L_\infty) \subseteq S_F$.

Finally, to see $\Miss(F) \subseteq S_F$: since $\pi: V \to Y$ is finite and dominant, $\pi$ is surjective. If $y \notin S_F = \pi(L_\infty)$, then $\pi^{-1}(y) \neq \emptyset$ and $\pi^{-1}(y) \cap L_\infty = \emptyset$, which implies $\pi^{-1}(y) \subset X$. Since $\pi$ is finite and the base point $y$ is a closed $\C$-point, every point in the fiber $\pi^{-1}(y)$ has residue field finite over $\C$, hence equal to $\C$ by the Nullstellensatz. Thus $F^{-1}(y) = \pi^{-1}(y) \neq \emptyset$, so $y \in F(X)$. Hence $\Miss(F) \subseteq S_F$.
\end{proof}

\begin{proposition}\label{prop:pure-codim-one}
$S_F = \emptyset$ or $S_F$ is a pure codimension-$1$ reduced hypersurface in $Y = \A^n_\C$.
\end{proposition}
\begin{proof}
$X = \A^n_\C$ is a dense affine open subscheme of the separated Noetherian scheme $V$. By Stacks Tag \href{https://stacks.math.columbia.edu/tag/0BCV}{0BCV} (Lemma 31.16.5), the complement of a dense affine open subscheme has all irreducible components of codimension 1; that is, for every generic point $\eta_E$ of an irreducible component $E \subset L_\infty$, $\dim \cO_{V, \eta_E} = 1 \implies \dim E = n-1$.

Since $\pi: V \to Y$ is finite, $\dim \pi(E) = \dim E = n-1$. If $L_\infty = \emptyset$, then $S_F = \emptyset$. Otherwise, by Proposition~\ref{prop:boundary-image}, $S_F = \pi(L_\infty) = \bigcup \pi(E_j)$ is a finite union of irreducible closed subvarieties of dimension $n-1$, hence a pure codimension-1 hypersurface in $\A^n_\C$.
\end{proof}

%=============================================================================
\section{Geometry over the Normal Locus}
%=============================================================================

Fix a point $y \in S_F^{\mathrm{nor}}$.

\begin{lemma}\label{lem:normal-neighborhood}
There exists a principal affine open neighborhood $U = D(g) \subset Y$ containing $y$ such that:
\begin{enumerate}[(1)]
\item $Z := S_F \cap U$ is normal, irreducible, and affine;
\item $X_U := F^{-1}(U) = D(g \circ F) \subset \A^n_\C$ is an affine variety;
\item $T := F^{-1}(Z) \subset X_U$ is a normal affine variety, and $T \to Z$ is an \'{e}tale morphism.
\end{enumerate}
\end{lemma}
\begin{proof}
Since $S_F$ is a reduced Noetherian scheme, the irreducible components of $S_F$ passing through $y$ correspond to the minimal prime ideals of the local ring $\cO_{S_F, y}$.
Since $\cO_{S_F, y}$ is a normal local ring, it is an integral domain, so exactly one irreducible component $S_0 \subset S_F$ passes through $y$.

Let $S' = \bigcup_{S_j \neq S_0} S_j$ be the union of all other irreducible components of $S_F$. Then $S'$ is closed in $S_F$ and $y \notin S'$.
The normal locus $S_F^{\mathrm{nor}}$ is open in $S_F$. Indeed, since $S_F$ is a reduced scheme of finite type over $\C$, its normalization
\[
\nu: S_F^\nu \longrightarrow S_F
\]
is finite.  Let
\[
\mathcal Q := \nu_* \cO_{S_F^\nu} / \cO_{S_F}.
\]
Then $\mathcal Q$ is coherent. For a point $x \in S_F$, one has $\mathcal Q_x = 0$ if and only if the normalization is an isomorphism in a neighborhood of $x$, equivalently if and only if $\cO_{S_F, x}$ is a normal local ring. Hence
\[
S_F^{\mathrm{nor}} = S_F \setminus \Supp(\mathcal Q),
\]
which is open because $\Supp(\mathcal Q)$ is closed.

Therefore, $\Omega := S_F^{\mathrm{nor}} \setminus S'$ is an open neighborhood of $y$ in $S_F$ with $\Omega \subset S_0$, so $\Omega$ is normal and irreducible.

Since $S_F \hookrightarrow Y$ is a closed immersion, there exists an open neighborhood $U_0 \subset Y$ containing $y$ such that $S_F \cap U_0 \subseteq \Omega$.
Since principal open subsets form a topological basis for $Y = \A^n_\C$, we can choose a principal affine open neighborhood $U = D(g) \subseteq U_0$ containing $y$.
Then $Z := S_F \cap U$ is normal, irreducible, and affine.

Since $U = D(g)$ is a principal open subset of $\A^n_\C$, $X_U = F^{-1}(U) = D(g \circ F)$ is a principal open affine subvariety of $\A^n_\C$. The preimage $T = F^{-1}(Z) = X_U \times_U Z$ is a closed subscheme of the affine variety $X_U$, hence $T$ is affine.

Furthermore, $F: X_U \to U$ is \'{e}tale, so the base change $T = X_U \times_U Z \to Z$ is \'{e}tale (Stacks Tag \href{https://stacks.math.columbia.edu/tag/02GH}{02GH}). Since $Z$ is normal and \'{e}tale extensions preserve normality (Stacks Tag \href{https://stacks.math.columbia.edu/tag/0336}{0336}, Lemma 10.163.9), $T$ is normal.
\end{proof}

Let $V_U := \pi^{-1}(U)$. Since $V$ is affine and $\pi$ is finite, $V_U$ is a normal affine variety.
We define the reduced closure of $T$ in $V_U$:
\begin{equation}\label{eq:closure-C}
C := \overline{T}^{V_U}_{\mathrm{red}} \subset V_U.
\end{equation}

\begin{lemma}\label{lem:closure-factorization}
The morphism $C \to U$ scheme-theoretically factors through $Z$, inducing a finite morphism $p: C \to Z$, and $T = C \cap X_U$ as schemes.
\end{lemma}
\begin{proof}
Since $T \subseteq \pi^{-1}(Z)$ and $\pi^{-1}(Z)$ is closed in $V_U$, the topological closure satisfies $|C| \subseteq |\pi^{-1}(Z)|$.
Let $\mathcal{I}_Z \subseteq \cO_U$ be the ideal sheaf of $Z$. The pullback of $\mathcal{I}_Z$ to $C$ is contained in every prime ideal of $C$. Since $C$ is reduced by definition, its nilradical is zero, so $\mathcal{I}_Z \cO_C = 0$. Thus $C \to U$ factors through $Z$, giving $p: C \to Z$. Since $\pi$ is finite, $p$ is finite.

For the scheme equality $T = C \cap X_U$: topologically, $|C \cap X_U| = \overline{|T|}^{V_U} \cap |X_U| = \overline{|T|}^{X_U} = |T|$ because $T$ is closed in $X_U$. Since both $T$ and $C \cap X_U$ are reduced subschemes of $X_U$ with the same underlying topological space, they are identical as schemes. Thus $C \setminus T = C \cap (V_U \setminus X_U)$.
\end{proof}

%=============================================================================
\section{The Boundary Collision Lemma}
%=============================================================================

\begin{theorem}[Boundary Collision Lemma]\label{thm:boundary-collision}
Let $a \in C \setminus T$ be a boundary point of $C$, and let $b = p(a) \in Z$. Then $b$ is a singular point of $Z$:
\[
p(C \setminus T) \subseteq \Sing(Z).
\]
\end{theorem}

\begin{proof}
Assume for contradiction that $b = p(a) \in Z_{\mathrm{reg}}$ is a smooth point of $Z$.

\medskip
\noindent\textbf{Step 1: Constructing the boundary prime divisor $E$.}
Since $a \in C \setminus T$ and $T = C \cap X_U$, we have $a \notin X_U$, so $a \in V_U \setminus X_U$.
The subvariety $X_U$ is a dense affine open subscheme of the separated Noetherian scheme $V_U$. By Stacks Tag \href{https://stacks.math.columbia.edu/tag/0BCV}{0BCV} (Lemma 31.16.5), every irreducible component of $V_U \setminus X_U$ passing through $a$ is a prime divisor in $V_U$. Let $E \subset V_U \setminus X_U$ be an irreducible component containing $a$. Then:
\[
\dim E = n-1, \qquad E \subset L_\infty \cap V_U, \qquad a \in E.
\]

\medskip
\noindent\textbf{Step 2: Constructing the interior closure prime divisor $R$.}
Since $a \in C = \overline{T}$, there exists an irreducible component $R \subseteq C$ containing $a$.
Since $T \subset C$ is a dense open subscheme, $R \cap T$ is a non-empty dense open subset of $R$, hence is irreducible.
We claim that $R \cap T$ is an irreducible component of $T$. Indeed, if $R \cap T \subsetneq Q \subseteq T$ for some irreducible closed subset $Q \subseteq T$, then taking closures in $C$ gives $R = \overline{R \cap T}^C \subseteq \overline{Q}^C \subseteq C$, which by maximality of $R$ implies $\overline{Q}^C = R$, so $Q \subseteq R \cap T$, a contradiction.

Since $T$ is normal and Noetherian, its irreducible components are pairwise disjoint open-and-closed normal integral schemes (Stacks Tag \href{https://stacks.math.columbia.edu/tag/0357}{0357}, Lemma 28.7.5). The restriction $R \cap T \to Z$ is \'{e}tale, so its image is a non-empty open subset of the irreducible variety $Z$, hence is dense. Thus $R \cap T \to Z$ is dominant and quasi-finite, which induces a finite algebraic field extension $K(Z) \subseteq K(R \cap T)$. Therefore:
\[
\dim R = \dim(R \cap T) = \dim Z = n-1 \implies \codim_{V_U} R = 1.
\]
Crucially, $R \cap X_U \neq \emptyset$ (since $R \cap T \neq \emptyset$), whereas $E \subset V_U \setminus X_U$ satisfies $E \cap X_U = \emptyset$. Therefore:
\[
R \neq E.
\]
Since $C \subseteq \pi^{-1}(Z)$ and $\pi(E) \subseteq \pi(L_\infty \cap V_U) \subseteq Z$, both $E$ and $R$ are contained in $\pi^{-1}(Z)$.

\medskip
\noindent\textbf{Step 3: Verification of Abhyankar Hypotheses.}
Since $b \in Z_{\mathrm{reg}}$, we may shrink $U$ to a smaller affine open neighborhood $W = \Spec A \subset U$ containing $b$ such that $D := Z \cap W = V(h)$ is a regular effective Cartier divisor in the regular scheme $W$, and $S_F \cap W = D$.

Let $V_W := \pi^{-1}(W)$ and let $E_W := E \cap V_W$, $R_W := R \cap V_W$. Since $a \in E \cap R \cap V_W$, $E_W$ and $R_W$ are non-empty irreducible prime divisors of $V_W$. If $E_W = R_W$, then taking closures in $V_U$ yields $E = \overline{E_W}^{V_U} = \overline{R_W}^{V_U} = R$, contradicting $E \neq R$. Thus $E_W \neq R_W$.

Moreover, since both $E$ and $R$ are contained in $\pi^{-1}(Z)$ and $V_W = \pi^{-1}(W)$, we have
\[
E_W, R_W \subseteq \pi^{-1}(Z) \cap \pi^{-1}(W) = \pi^{-1}(Z \cap W) = \pi^{-1}(D).
\]
Hence, setting $D' := (\pi^{-1}D)_{\mathrm{red}}$, both $E_W$ and $R_W$ are contained in the underlying topological space of $D'$.

Let $W^\circ := W \setminus D$. Then $W^\circ \cap S_F = \emptyset$.
By Proposition~\ref{prop:boundary-image}, $S_F = \pi(L_\infty)$, which implies $\pi^{-1}(W^\circ) \cap L_\infty = \emptyset$. Therefore,
\[
\pi^{-1}(W^\circ) = F^{-1}(W^\circ).
\]
Since $F$ is \'{e}tale and $\pi$ is finite, the restriction
\[
Y^\circ := \pi^{-1}(W^\circ) \longrightarrow W^\circ
\]
is a finite \'{e}tale covering.

The scheme $V_W = \Spec B$ is the normalization of $W$ in $K(X) = K(Y^\circ)$. Since $Y^\circ = \Spec B_h$, the integral closure of $A$ in $B_h$ is precisely $B$ (as $B$ is integrally closed in $K(X)$). Thus $V_W = \operatorname{Norm}_W(Y^\circ)$.

Since the residue field at every generic point of $D$ has characteristic 0, the corresponding finite separable extension of discretely valued fields is tamely ramified (Stacks Tag \href{https://stacks.math.columbia.edu/tag/0BSE}{0BSE} and Tag \href{https://stacks.math.columbia.edu/tag/09E9}{09E9}, Definition 15.113.7).

We now invoke \textbf{Stacks Lemma 58.31.6 (Tag \href{https://stacks.math.columbia.edu/tag/0EYH}{0EYH})}:
\begin{quote}
\emph{Let $W$ be a regular scheme, $D \subset W$ a regular effective Cartier divisor, and $Y^\circ \to W \setminus D$ a finite \'{e}tale morphism tamely ramified along $D$. Let $V_W = \operatorname{Norm}_W(Y^\circ)$. Then $D' := (\pi^{-1}D)_{\mathrm{red}}$ is an effective Cartier divisor on $V_W$, and $D'$ is a \textbf{regular scheme}.}
\end{quote}

In particular, for our point $a \in D'$, the local ring $\cO_{D', a}$ is a \textbf{regular local ring}, and therefore an \textbf{integral domain}.

\medskip
\noindent\textbf{Step 4: Minimal Primes and the Contradiction.}
By the preceding paragraph, $E_W, R_W \subseteq |D'|$. Both are irreducible codimension-1 closed subsets of $V_W$.
Since $D'$ is an effective Cartier divisor on the $n$-dimensional scheme $V_W$, every irreducible component of $D'$ has dimension $n-1$.
Specifically, $E_W \subseteq D_E \subseteq D'$ for some irreducible component $D_E$ of $D'$. Since $\dim E_W = \dim D_E = n-1$, we have $E_W = D_E$.
Similarly, $R_W = D_R$ for an irreducible component $D_R$ of $D'$.
Since $E_W \neq R_W$, $E_W$ and $R_W$ are two distinct irreducible components of the scheme $D'$.

Since $a \in E_W \cap R_W$, the point $a$ lies on two distinct irreducible components of $D'$. Consequently, the local ring $\cO_{D', a}$ has at least two distinct minimal prime ideals:
\[
\fp_E \neq \fp_R \subset \cO_{D', a}.
\]
Thus, $\cO_{D', a}$ has non-zero zero-divisors, and is \emph{not} an integral domain.

This directly contradicts Step 3, which established that $\cO_{D', a}$ is a regular local ring (hence a domain).

\medskip
We conclude that the assumption $b \in Z_{\mathrm{reg}}$ is false. Therefore, $b \in \Sing(Z)$, completing the proof.
\end{proof}

%=============================================================================
\section{Algebraic Hartogs Removal of the Boundary}
%=============================================================================

\begin{proposition}\label{prop:boundary-codim-two}
Let $\nu: \widetilde{C} \to C$ be the normalization of $C$. Then $\nu$ is finite, $T$ embeds as a dense open subscheme $T \subset \widetilde{C}$, and the boundary $B := \widetilde{C} \setminus T$ satisfies:
\[
\codim_{\widetilde{C}} B \ge 2.
\]
\end{proposition}
\begin{proof}
Since $C$ is a reduced affine scheme of finite type over $\C$, its normalization $\nu: \widetilde{C} \to C$ is a finite morphism (Stacks Tag \href{https://stacks.math.columbia.edu/tag/0BXQ}{0BXQ}). Since $T \subset C$ is normal (Lemma~\ref{lem:normal-neighborhood}), $\nu$ is an isomorphism over $T$ (Stacks Tag \href{https://stacks.math.columbia.edu/tag/0H7D}{0H7D}), so we identify $T$ with a dense open subscheme $T \subset \widetilde{C}$.

Let $g = p \circ \nu: \widetilde{C} \to Z$ be the composite finite morphism.
If $\tilde{a} \in B = \widetilde{C} \setminus T$, then $\nu(\tilde{a}) \in C \setminus T$. By Theorem~\ref{thm:boundary-collision},
\[
g(B) \subseteq \Sing(Z).
\]

Since $\widetilde{C}$ is a normal Noetherian scheme, its irreducible components $\widetilde{C} = \coprod_i \widetilde{C}_i$ are pairwise disjoint and open-and-closed (Stacks Tag \href{https://stacks.math.columbia.edu/tag/0357}{0357}, Lemma 28.7.5).
Since $T$ is dense in $C$, it meets every irreducible component of $C$. Under the finite normalization morphism $\nu: \widetilde{C} \to C$, each irreducible component of $\widetilde{C}$ dominates an irreducible component of $C$. Hence for each component $\widetilde{C}_i$, $T_i := T \cap \widetilde{C}_i$ is a non-empty dense open subset of $\widetilde{C}_i$. Since $T_i \to Z$ is \'{e}tale and $Z$ is irreducible, $g(T_i)$ is dense in $Z$. Since $g$ is finite, $g(\widetilde{C}_i) = Z$. Thus each $\widetilde{C}_i \to Z$ is finite dominant of dimension $\dim Z = n-1$.

Because $Z$ is normal, Serre's criterion $R_1$ (Stacks Tag \href{https://stacks.math.columbia.edu/tag/031S}{031S}) guarantees that all codimension-1 points of $Z$ are regular.
If $d = \dim Z \ge 2$, then $\codim_Z \Sing(Z) \ge 2 \implies \dim \Sing(Z) \le d - 2$.
For any irreducible component $B_0 \subseteq B \cap \widetilde{C}_i$, the restriction $B_0 \to g(B_0)$ is finite dominant, so $\dim B_0 = \dim g(B_0) \le \dim \Sing(Z) \le d - 2$.
Therefore, every irreducible component $B_0$ satisfies $\codim_{\widetilde{C}_i} B_0 = d - \dim B_0 \ge 2$, so $\codim_{\widetilde{C}_i}(B \cap \widetilde{C}_i) \ge 2$.
If $d \le 1$, then normality implies $\Sing(Z) = \emptyset$, so $B = \emptyset$.
\end{proof}

\begin{lemma}[Affine Hartogs Rigidity]\label{lem:affine-hartogs}
Let $X = \Spec A$ be an integral normal Noetherian affine scheme. Let $j: U \hookrightarrow X$ be a dense affine open subscheme such that $\codim_X(X \setminus U) \ge 2$.
Then $U = X$.
\end{lemma}
\begin{proof}
Every regular function $s \in \Gamma(U, \cO_U)$ defines a rational function in the fraction field $K = \Frac A$.
Since $\codim_X(X \setminus U) \ge 2$, every height-one prime ideal $\fp \subset A$ corresponds to a codimension-1 point in $U$.
Thus $s \in A_\fp$ for all $\operatorname{ht} \fp = 1$.
By the fundamental property of normal Noetherian domains (Stacks Tag \href{https://stacks.math.columbia.edu/tag/0AVB}{0AVB}, Lemma 15.24.18):
\[
A = \bigcap_{\operatorname{ht} \fp = 1} A_\fp \subset K.
\]
Thus $s \in A$, which proves $\Gamma(U, \cO_U) = A$.

Since $U = \Spec B$ is affine, the open immersion $j: U \hookrightarrow X$ corresponds to the ring homomorphism $A \to B = \Gamma(U, \cO_U)$, which is an isomorphism. By the anti-equivalence of affine schemes and commutative rings (Stacks Tag \href{https://stacks.math.columbia.edu/tag/01HX}{01HX}), $j: U \xrightarrow{\sim} X$ is an isomorphism of schemes. Thus $U = X$.
\end{proof}

\begin{theorem}\label{thm:closure-equality}
$T = C$. That is, $T$ is closed in $V_U$, and the restricted morphism $T \to Z$ is finite and \'{e}tale.
\end{theorem}
\begin{proof}
Since $\widetilde{C} \to Z$ is finite and $Z$ is affine, $\widetilde{C}$ is affine. Each disjoint component $\widetilde{C}_i$ is a normal Noetherian affine domain.
The scheme $T$ is affine (Lemma~\ref{lem:normal-neighborhood}), and $T_i = T \cap \widetilde{C}_i$ is an open-and-closed subscheme of the affine scheme $T$, hence $T_i$ is affine. (The affineness of $T_i$ is essential here.)

By Proposition~\ref{prop:boundary-codim-two}, $\codim_{\widetilde{C}_i}(\widetilde{C}_i \setminus T_i) \ge 2$. Applying Lemma~\ref{lem:affine-hartogs} to $T_i \hookrightarrow \widetilde{C}_i$, we obtain $T_i = \widetilde{C}_i$ for every component, so:
\[
T = \widetilde{C}.
\]

The normalization morphism $\nu: \widetilde{C} \to C$ is surjective and restricts to the open immersion $T \hookrightarrow C$.
Since $T = \widetilde{C}$, the image of $T \hookrightarrow C$ is the entire underlying space $|C|$.
An open immersion whose image is the entire topological space is an isomorphism. Thus:
\[
C = T.
\]
Therefore, $T$ is closed in $V_U$. Since $\pi: V_U \to U$ is finite, $\pi|_T: T \to Z$ is finite. Since $T \to Z$ was already \'{e}tale (Lemma~\ref{lem:normal-neighborhood}), $T \to Z$ is finite and \'{e}tale.
\end{proof}

%=============================================================================
\section{Surjectivity and the Conductor Principle}
%=============================================================================

\begin{proposition}[Non-Emptiness via Dominance]\label{prop:nonempty}
Let $S_i$ be the irreducible component of $S_F$ containing $Z$. Then $T = F^{-1}(Z) \neq \emptyset$.
\end{proposition}
\begin{proof}
Let $S_i = V(h)$, where $h \in \C[y_1, \dots, y_n]$ is an irreducible polynomial defining the hypersurface $S_i$. 

By Proposition~\ref{prop:keller-etale}, $F$ is dominant, so $F^*: \C[Y] \hookrightarrow \C[X]$ is injective. Therefore, $h \circ F \in \C[x_1, \dots, x_n]$ cannot be a constant. If $h \circ F = c \in \C$, then $F^*(h-c) = 0 \implies h = c$, contradicting that $h=0$ defines a hypersurface. 

Since $\Gamma(\A^n, \cO)^\times = \C^\times$, $h \circ F$ is a non-unit, so $F^{-1}(S_i) = V(h \circ F) \neq \emptyset$.

The restriction $F^{-1}(S_i) \to S_i$ is the base change of the \'{e}tale morphism $F$, hence is \'{e}tale. By Stacks Tag \href{https://stacks.math.columbia.edu/tag/02GH}{02GH} (Lemma 29.37.13), \'{e}tale morphisms are open, so the image
\[
I_i := F(F^{-1}(S_i)) \subseteq S_i
\]
is a non-empty open subset of $S_i$. Since $S_i$ is irreducible, $I_i$ is dense in $S_i$.

Since $Z \subseteq S_i$ is a non-empty open subset, the dense open set $I_i$ must intersect $Z$:
\[
I_i \cap Z \neq \emptyset \implies T = F^{-1}(Z) \neq \emptyset. \qedhere
\]
\end{proof}

\begin{theorem}[Surjectivity over $Z$]\label{thm:local-surjectivity}
The morphism $F^{-1}(Z) \to Z$ is surjective:
\[
F(F^{-1}(Z)) = Z.
\]
\end{theorem}
\begin{proof}
By Theorem~\ref{thm:closure-equality}, $T = F^{-1}(Z) \to Z$ is finite and \'{e}tale.
\begin{enumerate}[(1)]
\item A finite morphism is universally closed, so $F(T)$ is a closed subset of $Z$;
\item An \'{e}tale morphism is an open map, so $F(T)$ is an open subset of $Z$;
\item By Proposition~\ref{prop:nonempty}, $F(T) \neq \emptyset$.
\end{enumerate}
Since $Z$ is normal and irreducible, it is connected. In a connected topological space, the only non-empty clopen subset is the entire space. Therefore:
\[
F(T) = Z. \qedhere
\]
\end{proof}

\begin{proof}[Proof of Theorem~\ref{thm:normal-locus}]
Every point $y \in S_F^{\mathrm{nor}}$ has an affine open neighborhood $Z \subset S_F^{\mathrm{nor}}$ satisfying Lemma~\ref{lem:normal-neighborhood}. By Theorem~\ref{thm:closure-equality} and Theorem~\ref{thm:local-surjectivity}, $F^{-1}(Z) \to Z$ is finite, \'{e}tale, and surjective. Finiteness and \'{e}taleness are Zariski-local on the target.  Moreover, surjectivity follows pointwise: every point of $S_F^{\mathrm{nor}}$ lies in one of the open neighborhoods $Z$ over which the restricted morphism is surjective. Therefore the global restriction
\[
F^{-1}(S_F^{\mathrm{nor}}) \longrightarrow S_F^{\mathrm{nor}}
\]
is finite, \'{e}tale, and surjective.
\end{proof}

\begin{proof}[Proof of Corollary~\ref{cor:conductor}]
Let $y \in \Miss(F) = \A^n_\C \setminus F(\A^n_\C)$. By Proposition~\ref{prop:boundary-image}, $\Miss(F) \subseteq S_F$, so $y \in S_F$. 
By Theorem~\ref{thm:normal-locus}, every point of $S_F^{\mathrm{nor}}$ has at least one preimage in $\A^n_\C$. Since $y$ has no preimage, $y \notin S_F^{\mathrm{nor}}$. Therefore:
\[
y \in S_F \setminus S_F^{\mathrm{nor}} = \NonNormal(S_F).
\]

Let $\nu: S_F^\nu \to S_F$ be the normalization. The coherent sheaf quotient $\cQ = \nu_* \cO_{S_F^\nu} / \cO_{S_F}$ on stalks satisfies:
\[
0 \longrightarrow \cO_{S_F, y} \longrightarrow (\nu_* \cO_{S_F^\nu})_y \longrightarrow \cQ_y \longrightarrow 0.
\]
Since $\mathcal Q$ is coherent, $\mathcal Q_y=0$ if and only if there exists an open neighborhood $U\ni y$ such that $\mathcal Q|_U=0$.  On such a neighborhood,
\[
\mathcal O_U \cong (\nu_*\mathcal O_{S_F^\nu})|_U \cong (\nu|_{\nu^{-1}(U)})_* \mathcal O_{\nu^{-1}(U)}.
\]
The restriction $\nu^{-1}(U) \to U$ is finite, and a finite morphism is the relative spectrum of its pushforward algebra.  Hence the displayed isomorphism of $\mathcal O_U$-algebras implies that $\nu^{-1}(U)\to U$ is an isomorphism.

Since normalization commutes with restriction to open subsets, this is equivalent to $U$ being normal.  Thus
\[
\mathcal Q_y=0 \quad\Longleftrightarrow\quad \mathcal O_{S_F,y}\text{ is normal}.
\]
Consequently,
\[
\Supp(\mathcal Q)=\NonNormal(S_F).
\]
By Stacks Tag \href{https://stacks.math.columbia.edu/tag/00L2}{00L2} (Lemma 10.40.5), for any coherent module $\cQ$ on a Noetherian scheme, $\Supp(\cQ) = V(\Ann(\cQ)) = V(\fc_{S_F^\nu / S_F})$. Thus $\NonNormal(S_F) = V(\fc_{S_F^\nu / S_F})$.
\end{proof}

\begin{proof}[Proof of Corollary~\ref{cor:normal-surjectivity}]
If $S_F$ is globally normal, $S_F^{\mathrm{nor}} = S_F$. By Theorem~\ref{thm:normal-locus}, $F^{-1}(S_F) \to S_F$ is surjective, so $S_F \subseteq F(\A^n_\C)$. By Proposition~\ref{prop:boundary-image}, every point outside $S_F$ has a non-empty fiber. Hence $F(\A^n_\C) = \A^n_\C$.
\end{proof}

\begin{proof}[Proof of Corollary~\ref{cor:defect-rigidity}]
Let $W \subseteq S_F^{\mathrm{nor}}$ be a connected component. By Theorem~\ref{thm:normal-locus}, $F^{-1}(W) \to W$ is a finite \'{e}tale covering of non-empty connected base. The degree of a finite \'{e}tale cover over a connected base is constant: there exists $r_W \ge 1$ such that $\#F^{-1}(y) = r_W$ for all $y \in W$.
\end{proof}

%=============================================================================
\begin{thebibliography}{99}

\bibitem{BassConnellWright1982}
H.~Bass, E.~H.~Connell, and D.~Wright,
\emph{The Jacobian conjecture: reduction of degree and formal expansion of the inverse},
Bull. Amer. Math. Soc. (N.S.) \textbf{7} (1982), no.~2, 287--330.

\bibitem{BiaRosen1962}
A.~Bia{\l}ynicki-Birula and M.~Rosenlicht,
\emph{Injective morphisms of real algebraic varieties},
Proc. Amer. Math. Soc. \textbf{13} (1962), 200--203.

\bibitem{EGA4}
A.~Grothendieck,
\emph{\'{E}l\'{e}ments de g\'{e}om\'{e}trie alg\'{e}brique. IV. \'{E}tude locale des sch\'{e}mas et des morphismes de sch\'{e}mas. II},
Publ. Math. IH\'{E}S \textbf{24} (1965), 5--231.

\bibitem{Jelonek1993}
Z.~Jelonek,
\emph{The set of points at which a polynomial map is not proper},
Ann. Polon. Math. \textbf{58} (1993), no.~3, 259--266.

\bibitem{Jelonek1999}
Z.~Jelonek,
\emph{Testing sets for properness of polynomial mappings},
Math. Ann. \textbf{315} (1999), no.~1, 1--35.

\bibitem{Jelonek2022}
Z.~Jelonek,
\emph{A note on the Jacobian Conjecture},
Colloq. Math. \textbf{170} (2022), no.~1, 85--90.

\bibitem{JelonekLason2019}
Z.~Jelonek and M.~Laso\'{n},
\emph{Quantitative properties of the non-properness set of polynomial maps},
Adv. Math. \textbf{343} (2019), 743--766.

\bibitem{JelonekZoladek2020}
Z.~Jelonek and H.~\.{Z}o{\l}\k{a}dek,
\emph{A note on the Jacobian Conjecture},
arXiv:2004.05003 (2020, withdrawn).

\bibitem{Keller1939}
O.-H.~Keller,
\emph{Ganze Cremona-Transformationen},
Monatsh. Math. Phys. \textbf{47} (1939), 299--306.

\bibitem{StacksProject}
The Stacks Project Authors,
\emph{The Stacks Project},
\url{https://stacks.math.columbia.edu} (2026).
Tags cited: 
\href{https://stacks.math.columbia.edu/tag/02LR}{02LR} (Lemma 37.43.2),
\href{https://stacks.math.columbia.edu/tag/03GR}{03GR} (Lemma 29.54.15),
\href{https://stacks.math.columbia.edu/tag/0BXQ}{0BXQ} (Section 33.27),
\href{https://stacks.math.columbia.edu/tag/035F}{035F} (Lemma 29.54.1),
\href{https://stacks.math.columbia.edu/tag/02LS}{02LS} (Lemma 37.44.1),
\href{https://stacks.math.columbia.edu/tag/0BCV}{0BCV} (Lemma 31.16.5),
\href{https://stacks.math.columbia.edu/tag/02GH}{02GH} (Section 29.37),
\href{https://stacks.math.columbia.edu/tag/0336}{0336} (Lemma 10.163.9),
\href{https://stacks.math.columbia.edu/tag/0EYH}{0EYH} (Lemma 58.31.6),
\href{https://stacks.math.columbia.edu/tag/0BSE}{0BSE} (Section 58.31),
\href{https://stacks.math.columbia.edu/tag/09E9}{09E9} (Definition 15.113.7),
\href{https://stacks.math.columbia.edu/tag/0357}{0357} (Lemma 28.7.5),
\href{https://stacks.math.columbia.edu/tag/031S}{031S} (Lemma 10.157.4),
\href{https://stacks.math.columbia.edu/tag/0AVB}{0AVB} (Lemma 15.24.18),
\href{https://stacks.math.columbia.edu/tag/01HX}{01HX} (Section 26.6),
\href{https://stacks.math.columbia.edu/tag/0H7D}{0H7D} (Lemma 29.55.9),
\href{https://stacks.math.columbia.edu/tag/00L2}{00L2} (Lemma 10.40.5).

\end{thebibliography}

\end{document}
